\scanpage{89}
\mspage{44}
la dernière ligne dans (26) s'identifie aussi à
\[
(27)\qquad 1-\pi\longrightarrow E_{U_s}\longrightarrow E_{k(s)}-1
\]
associé à $U_s$, fibre de $U/S$ en $s$, sur le corps résiduel $k(s)$.

\uline{Résumons} La situation de la courbe relative $U/S$, et d'un $s\in S$ de codim $1$, définit par hensélisation (\inserted{\uline{\unsure{pas|non}}} stricte) en $s$ le diagramme (\unsure{26|25}), dont la dernière ligne est can. isom. à (27) (par \deleted{\unsure{l'inclusion}}\inserted{l'hom} \inserted{sur les $\pi_1$} correspondant à :
\[
\diagram{\begin{tikzcd}[ampersand replacement=\&]
U_s \arrow[r,no head] \arrow[d] \& U_{\widetilde O} \arrow[d] \\
\unsure{s|k} \arrow[r] \& \operatorname{Spec}\widetilde O
\end{tikzcd}}
\]
)

Une section rationnelle \inserted{$f_{\unsure{K}}$} de $U$ sur $S$, plus gén. $\widetilde f$ de $U_{\widetilde O}$ sur $\widetilde S$, définit une section \inserted{$E_{\widetilde f}$} de la ligne médiane de (26), ou ce qui revient au $\unsure{\overline m|\overset{m}{=}}$ un relèvement
\[
\diagram{\begin{tikzcd}[ampersand replacement=\&,column sep=large]
(28) \& E_{\widetilde K} \arrow[r,"\lambda"] \& E_{U_{\widetilde O}} \\
 \& \& E_{U_s} \arrow[u,"\sim"']
\end{tikzcd}}
\]
de $E_{\widetilde K}\longrightarrow E_{\widetilde O}$ ($\xleftarrow{\sim}E_{k(s)}$) --- où $E_{\widetilde K}$ peut d'ailleurs s'identifier au groupe de décomposition de $\unsure{\underline O_{S,s}|\underline U_{S,s}}$ dans $\overline K$. Les conditions suivantes sont
\margin{88}
